\chapter{When Things Go Wrong: Predictable Error Handling}

\section{Why Programs Crash}

\begin{kidbox}{Dropping a Glass of Milk}
Imagine you are carrying a glass of milk to the breakfast table, and your foot trips on a toy truck.
Crash! The glass hits the floor and milk spills everywhere!
What does a sensible person do?
Do you scream and burn down the entire house? Of course not!
You get a paper towel, wipe up the spilled milk, wash your hands, and pour a fresh glass!
In computer programs, when something goes wrong (like trying to read a file that does not exist or dividing by zero), bad languages burn down the house (they crash and quit!).
Shae teaches your computer how to wipe up the milk safely!
\end{kidbox}

\section{The Banishment of Null}

\begin{kidbox}{The Empty Birthday Present Box}
Imagine it is your birthday. A friend hands you a beautifully wrapped box with a golden bow.
You smile excitedly and unwrap the paper. You open the lid...
And inside there is nothing. Not even air! It is a black hole that makes your bedroom vanish!
In 1965, a computer scientist named Sir Tony Hoare invented \texttt{null}. Decades later, he called it his \textbf{"Billion-Dollar Mistake"} because unhandled null pointers have crashed millions of computers and caused billions of dollars in software bugs.
Shae banishes raw, dangerous \texttt{null} pointers! Instead, Shae gives you explicit \texttt{Option} and \texttt{Result} boxes!
\end{kidbox}

\section{Structured Exceptions: Try, Catch, and Throw}

For unexpected runtime conditions, Shae provides structured exception handling:

\begin{lstlisting}[language=shaelang]
fn read_secure_file(path) {
    if path == "" {
        throw "Path cannot be empty!";
    }
    // Attempt file reading...
    return "Secret file contents";
}

try {
    let data = read_secure_file("");
    println(data);
} catch (error) {
    println("Caught an error gracefully: " + error);
}

println("Program continues running safely!");
\end{lstlisting}

\begin{deepbox}{Stack Unwinding and Exception Frames}
When a \texttt{throw} statement executes:
\begin{enumerate}
    \item The runtime captures the thrown value and constructs an error diagnostic backtrace.
    \item In the Bytecode VM, the engine inspects the current \texttt{CallFrame}. If no active exception handler is registered in the frame, the frame is popped, local variables are dropped, and execution unwinds to the caller frame.
    \item Unwinding continues recursively until an active handler frame is encountered.
    \item The instruction pointer \texttt{ip} is jumped to the catch block offset, the caught error value is pushed onto the VM stack, and normal execution resumes.
    \item If the call stack unwinds completely without finding any catch handler, a clean fatal error diagnostic is printed to \texttt{stderr} with full source file, line, and column spans.
\end{enumerate}
\end{deepbox}
